\chapter{Domaines, codomaines et hypothèses des résultats}
\label{app:premisas}

\begin{enumerate}
\item Le recensement \(104\,976\to468\to243\), l'image de 42 triplets, la
factorisation \(468=26\times18\), les cinq régions et les cinq cycles sont des
théorèmes exhaustifs du système fini de transport spécifié dans cet
ouvrage.

\item Le catalogue enrichi produit, par des prédicats finis de fibre et sans
consultation décimale, l'unique orbite de clôture, l'unique orbite de
propagation et l'unique orbite d'auto-échelle ; le calibre interne y oriente
les représentants \(010211,201101,121200\). Leurs caractères sont
construits, respectivement, par la clôture pentarégionale avec retour
unitaire, la récurrence de semi-groupe \((n+1)a_{n+1}=a_n\) et l'unique rayon positif
de l'incidence \(F_{\rm av}\). L'évaluation archimédienne publie
alors les valeurs engendrées et la reconnaissance conventionnelle les nomme
\(\pi,e,\varphi\). Pour chaque profondeur \(N\), le prolongement TPK produit
un préfixe unique compatible avec celui de profondeur \(N-1\) ; la limite
inverse de ces préfixes fixe le développement complet. Ainsi, les mots,
les caractères et les valeurs sont des étapes successives d'une construction
interne, non des étiquettes choisies au moyen de tables externes.

\item Le théorème dodécaphasique a pour domaine causal l'état terminal
enrichi \(x_{\rm term}\) produit par APP--TRIT--TPK. La carte signée
publie à partir de lui
\(x_{\rm term}\xrightarrow{R_{\rm sgn}}U_{12}^{\rm sgn}\), et l'inversion
entière publie ensuite
\(U_{12}^{\rm sgn}\xrightarrow{\mathcal H_{12}^{-1}}K\). Dans ce domaine,
les correspondances avec \((D_3K,D_4K,Q)\), la représentation de Hadamard et
la clôture dodécaphasique sont exactes et réversibles. Ni \(U\) ni \(K\) ne sont des
entrées génératrices. La trace réduite d'une trajectoire positionnelle et le
système bivarié complet sont des objets distincts, comme l'établit le
\cref{chap:tpk-definicion}.

\item L'égalité avec la chaîne centrale 2018 est une reconnaissance externe.
Le mot mathématique interne et sa démonstration sont indépendants des
recommandations métrologiques ultérieures.

\item L'instance finie de la double projection démontre un drapeau commun,
\(W_{12}\) et \(M_{12}\). Dans la branche réticulaire, on fixe le recollement de
\(A_2^{12}\), la chambre radiale et l'origine ; le 3-voisin résultant est certifié
pair, unimodulaire et sans racines, et est identifié classiquement à Leech. La
construction de \(V^\natural\) et du Monstre est une conséquence classique. La
flèche \(W_{12}\to W_{24}\) est réalisée au moyen du plan résiduel d'une
dodécade et du relèvement unique de ses hexades, relativement à un choix
\((D,\ell)\) ; les identifications s'effectuent au moyen d'applications, non
par des égalités de cardinalité. La compatibilité projective de ces
données et leur double réalisation sur l'état enrichi commun sont démontrées
dans les
\cref{p12-83-c17-s2:thm:lector-excepcional-global-rev6,p12-83-c17-s2:thm:doble-realizacion-global-rev6}.

\item Le théorème des cylindres produit un unique réel pour chaque histoire de son
domaine. La loi de prolongement couvre de manière cohérente chaque compact et
l'exhaustion dirigée couvre la droite entière, comme le démontrent les
\cref{p12-83-c18-s1:thm:sobreyectividad-evaluacion-compacto,p12-83-c18-s1:prop:agotamiento-recta-real-u024,p12-83-c18-s2:hmtch:thm-cociente-real}.
Par conséquent, l'évaluation est surjective et son quotient archimédien est
\(\mathbb R\) ; la couverture est une conclusion du générateur HMT, non une
prémisse supplémentaire.

\item Le produit natif est démontré dans le secteur canonique des mots.
L'extension monoïdale à l'espace complet des histoires et l'entrelacement
avec les filtrations de survie sont formulés comme des applications distinctes
dans le \cref{chap:producto-nativo}.
\end{enumerate}
